\subsection{External Benchmarks and Multilingual Encoders}
\label{sec:external-benchmark-plan}

To test whether domain-pretrained initialization helps consistently across recruitment tasks, the experiments compare XLM-R-large~\citep{conneau2020xlmr} and ESCOXLM-R~\citep{zhang2023escoxlmr}. They use six public datasets: SkillSpan~\citep{zhang2022skillspan}, Kompetencer~\citep{jensen2022kompetencer}, German ICT~\citep{gnehm2022transfer}, Green~\citep{green2022benchmark}, Sayfullina~\citep{sayfullina2018softskills}, and FIJO~\citep{beauchemin2022fijo}. Separate skill and knowledge streams in SkillSpan and Kompetencer give eight extraction tasks, each retaining its original entity types.

Both encoders use fresh linear first-subword BIO heads, AdamW with a learning rate of $2\times10^{-5}$, batch size 8, and a maximum of 20 epochs. Training uses seeds 42, 43, and 44, with development typed exact token-span micro-F1 selecting checkpoints. For the Chinese comparison, both encoders and their new nine-label BIO heads were fine-tuned on B2's 2,150 training and 169 development records. Both models process the original Chinese characters directly. The new BIO head labels the first subword of each character; predictions from chunks of at most 128 characters are joined, including the remainder of long sentences. Tokenizer placeholders preserve positions for whitespace and otherwise discarded characters; scoring uses the original text and offsets. Development F1 selects the checkpoint. Gold150 is used only for diagnostic scoring within these runs, although it was previously used in guideline development. Exact text and identifier checks found no overlap across the three partitions; advertisement-level independence is not established.

The Chinese comparison uses the same optimizer settings, maximum epochs, and three seeds as above. Its linear heads and training budget differ from JobBERT's CRF and Qwen LoRA, so shared-reference scores support descriptive comparisons rather than a controlled ranking of those architectures. We report typed exact and relaxed character-span F1, boundary-only F1, and macro-F1 over L/K/S/T. L has only two reference spans but receives one quarter of macro-F1 weight.

External-task paired sentence-bootstrap intervals estimate the mean ESCOXLM-R minus XLM-R difference for three fixed trained seeds (2,000 resamples; seed 20260921). We report training-seed SD separately. Intervals omit advertisement clustering and multiple-comparison adjustment.

\hyperref[app:e]{Appendix E} gives separate protocols for five-seed MaChAmp-CRF reruns on SkillSpan and German ICT, five-seed Kompetencer classification, and single-seed NNOSE retrieval. These preserve their task-specific scoring: Danish BIO extraction differs from fine-grained classification of supplied spans, and the public German ICT task~\citep{gnehm2022transfer} is distinct from the EDU/EXP/LNG extraction and taxonomy-classification pipeline~\citep{gnehm2022skills}. Languages and annotation schemes differ, so absolute scores do not rank extraction difficulty.

The Chinese encoder comparison uses 10,000 draws (seed 20260924), jointly resampling the same 150 sentences for both models and all three fixed seeds. Each draw averages paired differences in recomputed per-seed micro-F1. This post hoc interval describes reference-sentence sampling uncertainty, not population generalization, and omits advertisement clustering.
